\documentclass[11pt,a4paper,twoside]{article}
\usepackage[utf8]{inputenc}
\usepackage[english]{babel}
\usepackage{actuarialangle, actuarialsymbol}
\usepackage{amsmath,amsfonts,amsthm,amssymb,mathtools}

\makeatletter
\newcommand{\ix}{\@actothersc{i}}
\makeatother

\begin{document}

In life insurance what are rights and obligations of the parties on contractual alterations?


Contractual Alterations by insurer:
Policies in life insurance typically have long policy terms. During the contract duration,
the insurer can adjust premiums, benefits, or policy conditions only in a very restricted
way, see VVG, § 163
The insurance company is allowed to recalculate the premium, if
1. the required benefits have changed permanently and in an unforeseeable way compared to
the calculation bases used for the agreed upon premiums,
2. the newly calculated (via new calculation bases) premium is appropriate and needed in
order to make sure the agreed upon benefits can be paid for
3. and an independent trustee verified 1., 2.

In practice, the most common examples of contractual alterations are:
Repurchase: cancellation of the contract
Release from premium payment: no future premiums, reduction of benefits
- Change of the benefits
- Change of policy term: premium payment period or policy term are changed (e.g. shortened)

If a contract in life insurance is cancelled prematurely, the insured person has a claim on a
certain amount of money, called the surrender value, see VVG, § 169. The process itself is called repurchase.

In general, for a term life insurance, there is no surrender value (this would otherwise lead to selective cancellation)

According to the law, the surrender value is given by the reserve; in case of cancellation
it is at least the reserve that would result from a distribution of the acquisition costs on the
first 5 years upon conclusion of contract.


Alterations by insuree:
By law, a policy holder can demand that their policy is transformed into a policy without
premiums until the end of the current policy period (also called release from premia).
In this case, they must accept reduced benefits or reduced insurance protection.
The reduced benefits are calculated according to the principle of equivalence with the
help of the surrender value as a single premium (VVG, § 165).



Give examples for computing the surrender value and benefits after release of premia.

Consider a mixed capital insurance for a \( 50 \) year old man:
Policy term: \( 15 \) years
Insured sum: \( 50 000 \)
with annual premiums in advance during the whole policy term.
Repurchase after 10 years.


Surrender value for one-time acquisition costs:
The annual sufficient premium is determined as before (here the commutation terms are already inserted)
\[
  \actsymb[][a]{P}{50} = \frac{50000}{(1 - \beta)(N_{50} - N_{65} - \alpha^Z \cdot 15 \cdot D_{50})}(D_{65} + M_{50} - M_{65} + (\alpha^\gamma + \gamma_1)(N_{50} - N_{65}) 
\]

and the sufficient reserve for ten years is

\[
  \Vx[a][10]{50} = \frac{1}{D_{60}}\left(50000 \left(M_{60} - M_{65} + D_{65} + (\alpha^\gamma + \gamma_1)(N_{60} - N_{65}) -   \actsymb[][a]{P}{50}(1 - \beta)(N_{60} - N_{65} \right) \right)
\]


Surrender value with acquisition costs distributed on 5 years: \\

The premia can be found by solving the following equation for \( \actsymb[][a]{P}{50} \) and multiplying the result with \( 50000 \)
\[
  \actsymb[][a]{P}{50} \ax**{50:\angl 15} = \Ax{50:\angl 15} + \ax**{50:\angl 15} \left(\alpha^Z \frac{15 \cdot \actsymb[][a]{P}{50}}{5}  + \beta \actsymb[][a]{P}{50} + (\alpha^\gamma + \gamma_1) \right)
\]
Inserting this into the previous equation for \(   \Vx[a][10]{50} \) then gives the surrender value. It will typically be higher than the method one-time acquisition costs
for the first five years.


Consider again the previous setup, but now after \( 10 \) years the policyholder requests a release of premia, which would give a mixed
capital insurance with a policy term of \( 5 \) years and a reduced insured sum of \( \tilde{S} \).

In order to calculate \( \tilde{S} \), we consider the surrender value after \( 10 \) years and the sufficient
reserve (with acquisition costs distributed on the first \( 5 \) years).
This amount is considered as single premium (acquisition costs and collection costs are to
be neglected) for the transformed insurance.

\[ 
  \Vx[a][10]{50} = \tilde{S}( \Ax{60:\angl 5} + (\alpha^\gamma + \gamma_1) \ax**{60:\angl 5}) \approx \tilde{S} \cdot 0.93059
\]
solving for \( \tilde{S} \) then gives the remaining benefits.


How is the regulation for unisex tariffication included in the reserve building?

The realization of the principle of equal treatment of men and women in relation to goods
and services is the content of the directive 2004/113/EG.

Potential Solutions:
1. Continue as before!
The reserve is simply calculated with gender-specific tables
Problem: Different calculation bases for tariffing and reservation
2. Careful tariffing and reservation
The table is chosen which results in higher premiums and reserves
3. Unisex reservation:
For both tariffing and reservation, unisex tables are chosen.
When constructing the tables, the temporal development of the mix must be considered.


What are the sources and treatment of surpluses in insurance companies?

Sources are often a consequence due to the especially careful choice of calculation bases:
- Interest Surplus: The interest rate realized at the capital market is higher than the
technical interest rate (Zinsergebnis)
- Risk Surplus: The real, observed mortality rates are higher/lower than those used for the
calculation of the products (Risikoergebnis)
- Other Surplus: The costs used for the calculation are higher than the real costs

The sum of the different surplus components is the gross surplus. It can be used in terms of profit participation follows:
- Formation of equity capital
- Payment of a dividend to the shareholders (owners of the company)
- Feed the RfB (Rückstellung für Beitragsrückerstattung; reserve for profit participation for policy holders)
- Participation of the policy holders via direct credit (reduced premium), interest-bearing accrual/investment, direct payment, increased benefits

By legislation if there's a gross surplus there's at least a minimal participation of the insurees.

The actuary responsible makes a proposal to the board, and the board then decides the surplus
declaration for the following year.
\end{document}

